% Page 045

\seite{45}

\abschnitt{Ernte}
\pstart
In diessiger Gegend kann es allein als noch befriedigend angesehen\ob
werden. Jedoch auch die Preise der erste und Roggen aufs\ob
Drittel und Grünfutter konnten wegen des Mangels nicht als in\ob
frühern Summarisirt werden, ein grosser Teil wurde davon\ob
noch Geraden gebracht worden. Wegen vielfach ungünst.\ob
Witterung zog sich die Grummeternte in die Länge. Obst gab es\ob
ziemlich wenig auf eine, was fast die Hälfte verfaulte ganz das\ob
Bedürfnis fiel im allgemeinen nicht. Die Kartoffeln waren\ob
bei einem mittleren Ertrag.

Gelesen 14.\ 12.\ 08

\pend
\begin{verse}
1908
\end{verse}
\pstart

Weihnachtsfeier.\ob
Am Geburtstag Sr. Maj. des Kaisers wurden den\ob
nächst Besuch der Gemein eingeladenen Bürgern zu\ob
Seffern mithin Kinder des Herrn Ortsschulinspektor Pfarrer\ob
geziemte Aufführsten mittheilen, die von den grössern\ob
Kindern einstudirten, etlichen Gedichte, die von einigen\ob
Kindern vorgetragen wurden. Die Ansprachen hielt\ob
Kindergeld. Lehrer Kespen aus Schleid.

Prüfung.\ob
Die diesjähr. Prüfe wurde am 14.\ Dezember 08 durch\ob
den Herrn Kreisschulinspektor Herz aus Bitburg erwähnt.

Gelesen 30/3.\ 09
Schmitt, Pfr.

Osterprüfung.\ob
Die Osterprüfung wurde am 30.\ März durch den Herrn Orts\ohb
schulinspektor Herrn Pfarrer Schmitt aus Seffern abgehalten.

Das Schuljahr 1909--10.\ob
Entlassen wurden in diesem Jahre 5 Knaben. Ausge\ohb
den widmen sich der Landwirthschaft. Aufgenommen war
\pend
